\chapter{Translations et vecteurs}

\noindent\textit{Un vecteur code un \textbf{déplacement} : une direction, un sens et une
longueur. Il permet de décrire une translation, d'additionner des déplacements, et de
démontrer des propriétés de figures (parallélogrammes) sans mesurer d'angle. Les
coordonnées d'un vecteur dans un repère en font aussi un outil de calcul très efficace.}
\vspace{8pt}

\connecteBanner

\section{Vecteurs : direction, sens, norme}

\begin{definition}[Vecteur]
Le \textbf{vecteur} $\vec{AB}$ représente le déplacement qui amène le point $A$ sur le
point $B$. Il est caractérisé par trois éléments : sa \textbf{direction} (la droite $(AB)$),
son \textbf{sens} (de $A$ vers $B$) et sa \textbf{norme} (la longueur $AB$, notée
$\|\vec{AB}\|$).
\end{definition}

\begin{illus}
\begin{tikzpicture}[>=Stealth]
\draw[onbo,->,line width=1.4pt] (0,0) -- (3,1.5);
\node[below] at (0,0) {$A$};
\node[above] at (3,1.5) {$B$};
\draw[onbblue,->,line width=1.4pt] (1,-2.5) -- (4,-1);
\node[below] at (1,-2.5) {$C$};
\node[above] at (4,-1) {$D$};
\node[onbo,above] at (1.4,1) {\small $\vec{AB}$};
\node[onbblue,above] at (2.3,-1.9) {\small $\vec{CD}$};
\end{tikzpicture}
\fig{Figure — $\vec{AB}$ et $\vec{CD}$ ont même direction, sens et longueur : même
vecteur.}
\end{illus}

\begin{remarque}
Deux vecteurs $\vec{AB}$ et $\vec{CD}$ sont \textbf{égaux} s'ils ont même direction, même
sens et même longueur, même si $AB$ et $CD$ ne sont pas au même endroit dans le plan.
\end{remarque}

\section{Translation}

\begin{definition}[Translation]
La \textbf{translation de vecteur $\vec u$} associe à tout point $M$ le point $M'$ tel que
$\vec{MM'}=\vec u$. Le point $M'$ est l'\textbf{image} de $M$ par cette translation.
\end{definition}

\begin{illus}
\begin{tikzpicture}[>=Stealth]
\coordinate (A) at (0,0);
\coordinate (B) at (2,0);
\coordinate (C) at (0.5,1.8);
\coordinate (Ap) at (3,-1);
\coordinate (Bp) at (5,-1);
\coordinate (Cp) at (3.5,0.8);
\draw[onbo,line width=1.2pt] (A) -- (B) -- (C) -- cycle;
\draw[onbblue,line width=1.2pt] (Ap) -- (Bp) -- (Cp) -- cycle;
\draw[onbmuted,->,dashed,line width=0.9pt] (A) -- (Ap);
\draw[onbmuted,->,dashed,line width=0.9pt] (B) -- (Bp);
\draw[onbmuted,->,dashed,line width=0.9pt] (C) -- (Cp);
\node[below left] at (A) {$A$};
\node[below right] at (B) {$B$};
\node[above] at (C) {$C$};
\node[below] at (Ap) {$A'$};
\node[below right] at (Bp) {$B'$};
\node[above] at (Cp) {$C'$};
\node[onbmuted,above] at (1.5,-0.75) {\small $\vec u$};
\end{tikzpicture}
\fig{Figure — Translation de $ABC$ par $\vec u$ : $\vec{AA'}=\vec{BB'}=\vec{CC'}=\vec u$.}
\end{illus}

\begin{propriete}[Propriétés conservées]
Une translation conserve les longueurs, les angles, le parallélisme et les aires : l'image
d'une figure par translation est une figure \textbf{superposable} à la figure de départ.
\end{propriete}

\section{Somme de deux vecteurs, relation de Chasles}

\begin{propriete}[Relation de Chasles]
Pour tous points $A$, $B$, $C$ :
\[
\vec{AB}+\vec{BC}=\vec{AC}.
\]
\end{propriete}

\begin{illus}
\begin{tikzpicture}[>=Stealth]
\coordinate (A) at (0,0);
\coordinate (B) at (3,2);
\coordinate (C) at (5,6);
\draw[onbo,->,line width=1.3pt] (A) -- (B);
\draw[onbblue,->,line width=1.3pt] (B) -- (C);
\draw[onbgreen,->,dashed,line width=1.3pt] (A) -- (C);
\node[below] at (A) {$A$};
\node[right] at (B) {$B$};
\node[above] at (C) {$C$};
\node[onbo,below right] at (1.5,1) {\small $\vec{AB}$};
\node[onbblue,right] at (4,4) {\small $\vec{BC}$};
\node[onbgreen,left] at (2,3.3) {\small $\vec{AC}$};
\end{tikzpicture}
\fig{Figure — Relation de Chasles : $\vec{AB}+\vec{BC}=\vec{AC}$.}
\end{illus}

\begin{remarque}
La relation de Chasles se généralise à plusieurs points : par exemple
$\vec{AB}+\vec{BC}+\vec{CD}=\vec{AD}$. On l'utilise aussi pour simplifier une somme :
$\vec{AB}+\vec{BA}=\vec{AA}=\vec 0$ (le \textbf{vecteur nul}).
\end{remarque}

\section{Coordonnées d'un vecteur}

\begin{propriete}[Coordonnées d'un vecteur]
Dans un repère, si $A(x_A\,;y_A)$ et $B(x_B\,;y_B)$, le vecteur $\vec{AB}$ a pour
coordonnées
\[
\vec{AB}\begin{pmatrix}x_B-x_A\\y_B-y_A\end{pmatrix}, \qquad
\|\vec{AB}\|=\sqrt{(x_B-x_A)^{2}+(y_B-y_A)^{2}}.
\]
Pour additionner deux vecteurs, on additionne leurs coordonnées, terme à terme.
\end{propriete}

\begin{illus}
\begin{tikzpicture}
\begin{axis}[width=7.5cm,height=6cm,axis lines=middle,xlabel={\small $x$},ylabel={\small $y$},
  xmin=-1,xmax=7,ymin=-1,ymax=8,xtick={2,6},ytick={3,7}]
\draw[onbmuted,dashed] (axis cs:2,3) -- (axis cs:6,3) -- (axis cs:6,7);
\draw[onbo,->,line width=1.3pt] (axis cs:2,3) -- (axis cs:6,7);
\node[below] at (axis cs:4,3) {\small $4$};
\node[right] at (axis cs:6,5) {\small $4$};
\end{axis}
\end{tikzpicture}
\fig{Figure — Coordonnées d'un vecteur : $\vec{AB}$ a pour coordonnées $(4\,;4)$.}
\end{illus}

\begin{exemple}
$A(1\,;2)$, $B(4\,;6)$ : $\vec{AB}\begin{pmatrix}4-1\\6-2\end{pmatrix}=\vec{AB}\begin{pmatrix}3\\4\end{pmatrix}$,
et $\|\vec{AB}\|=\sqrt{3^{2}+4^{2}}=5$.
\end{exemple}

\section{Démontrer qu'un quadrilatère est un parallélogramme}

\begin{propriete}[Parallélogramme et vecteurs]
Un quadrilatère $ABCD$ est un parallélogramme si et seulement si $\vec{AB}=\vec{DC}$
(côtés $[AB]$ et $[DC]$ parallèles et de même longueur, dans le même sens).
\end{propriete}

\begin{exemple}
$A(0\,;0)$, $B(5\,;2)$, $C(8\,;6)$, $D(3\,;4)$ : $\vec{AB}\begin{pmatrix}5\\2\end{pmatrix}$
et $\vec{DC}\begin{pmatrix}8-3\\6-4\end{pmatrix}=\begin{pmatrix}5\\2\end{pmatrix}$. Les deux
vecteurs sont égaux, donc $ABCD$ est un parallélogramme.
\end{exemple}

% ═══════════════════════════════════════════════════════════════
\section{L'essentiel à retenir}

\begin{synthese}
\textbf{Vecteur.} $\vec{AB}$ = direction, sens, norme ($\|\vec{AB}\|=AB$).\\[4pt]
\textbf{Translation.} $\vec{MM'}=\vec u$ pour tout point $M$ ; conserve longueurs, angles,
aires, parallélisme.\\[4pt]
\textbf{Chasles.} $\vec{AB}+\vec{BC}=\vec{AC}$ (se généralise à plusieurs points).\\[4pt]
\textbf{Coordonnées.} $\vec{AB}\begin{pmatrix}x_B-x_A\\y_B-y_A\end{pmatrix}$ ;
$\|\vec{AB}\|=\sqrt{(x_B-x_A)^{2}+(y_B-y_A)^{2}}$ ; addition coordonnée par coordonnée.\\[4pt]
\textbf{Parallélogramme.} $ABCD$ parallélogramme $\iff\vec{AB}=\vec{DC}$.\\[4pt]
\textbf{Piège fréquent.} $\vec{AB}\ne\vec{BA}$ (sens opposé) ; pour un parallélogramme,
bien respecter l'ordre des sommets dans $\vec{AB}=\vec{DC}$ (pas $\vec{CD}$).
\end{synthese}

% ═══════════════════════════════════════════════════════════════
\section{Méthodes \& savoir-faire}

\methode{Reconnaître si deux vecteurs sont égaux}
Vérifier qu'ils ont même direction (parallèles), même sens et même longueur — ou, en
coordonnées, les mêmes coordonnées.
\begin{exemple}
$\vec{AB}\begin{pmatrix}2\\5\end{pmatrix}$ et $\vec{CD}\begin{pmatrix}2\\5\end{pmatrix}$ :
mêmes coordonnées, donc $\vec{AB}=\vec{CD}$.
\end{exemple}

\methode{Construire l'image d'un point par une translation}
Reporter le vecteur $\vec u$ à partir du point, en respectant sa direction, son sens et sa
longueur (ou additionner ses coordonnées).
\begin{exemple}
$A(2\,;5)$, $\vec u\begin{pmatrix}3\\-4\end{pmatrix}$ : $A'(2+3\,;5-4)=A'(5\,;1)$.
\end{exemple}

\methode{Calculer les coordonnées d'un vecteur}
Soustraire les coordonnées de l'origine à celles de l'extrémité.
\begin{exemple}
$A(-2\,;1)$, $B(3\,;13)$ : $\vec{AB}\begin{pmatrix}3-(-2)\\13-1\end{pmatrix}=\begin{pmatrix}5\\12\end{pmatrix}$.
\end{exemple}

\methode{Calculer la norme d'un vecteur}
Appliquer $\|\vec{AB}\|=\sqrt{(x_B-x_A)^{2}+(y_B-y_A)^{2}}$ (théorème de Pythagore).
\begin{exemple}
$\vec{AB}\begin{pmatrix}5\\12\end{pmatrix}$ : $\|\vec{AB}\|=\sqrt{25+144}=\sqrt{169}=13$.
\end{exemple}

\methode{Additionner deux vecteurs}
Additionner leurs coordonnées terme à terme, ou utiliser la relation de Chasles si les
vecteurs sont donnés par des points.
\begin{exemple}
$\vec u\begin{pmatrix}3\\-2\end{pmatrix}$, $\vec v\begin{pmatrix}1\\5\end{pmatrix}$ :
$\vec u+\vec v\begin{pmatrix}4\\3\end{pmatrix}$.
\end{exemple}

\methode{Utiliser la relation de Chasles pour simplifier une somme}
Repérer les points qui se répètent consécutivement (fin d'un vecteur = début du suivant) et
les « annuler ».
\begin{exemple}
$\vec{AB}+\vec{BC}+\vec{CA}=\vec{AA}=\vec 0$.
\end{exemple}

\methode{Démontrer qu'un quadrilatère est un parallélogramme}
Calculer les coordonnées de $\vec{AB}$ et de $\vec{DC}$ ; si elles sont égales, conclure que
$ABCD$ est un parallélogramme.
\begin{exemple}
$A(1\,;1)$, $B(6\,;2)$, $C(9\,;6)$, $D(4\,;5)$ : $\vec{AB}\begin{pmatrix}5\\1\end{pmatrix}$,
$\vec{DC}\begin{pmatrix}5\\1\end{pmatrix}$. Égaux : $ABCD$ est un parallélogramme.
\end{exemple}

% ═══════════════════════════════════════════════════════════════
\section{Exercices d'application}
\begin{multicols}{2}

\rubrique{Vecteurs et translations}
\exo{}\diff{1} $\vec{AB}\begin{pmatrix}2\\5\end{pmatrix}$ et
$\vec{CD}\begin{pmatrix}2\\5\end{pmatrix}$. Que peut-on dire de ces deux vecteurs ?

\exo{}\diff{2} Calculer les coordonnées de l'image $A'$ du point $A(2\,;5)$ par la
translation de vecteur $\vec u\begin{pmatrix}3\\-4\end{pmatrix}$.

\exo{}\diff{2} Le triangle $ABC$ a pour sommets $A(0\,;0)$, $B(4\,;1)$, $C(2\,;5)$. On le
translate par le vecteur $\vec u\begin{pmatrix}3\\-2\end{pmatrix}$. Donner les coordonnées
de $A'$, $B'$, $C'$.

\exo{}\diff{2} Une translation transforme $E$ en $E'(7\,;-1)$. Sachant que le vecteur de la
translation est $\vec u\begin{pmatrix}-2\\5\end{pmatrix}$, calculer les coordonnées de $E$.

\rubrique{Coordonnées et norme d'un vecteur}
\exo{}\diff{1} $A(1\,;2)$, $B(4\,;6)$. Calculer les coordonnées de $\vec{AB}$ et sa norme.

\exo{}\diff{2} $A(-2\,;1)$, $B(3\,;13)$. Calculer les coordonnées de $\vec{AB}$ et sa
norme.

\exo{}\diff{2} $A(0\,;0)$, $B(-6\,;8)$. Calculer les coordonnées de $\vec{AB}$ et sa norme.

\exo{}\diff{2} $A(2\,;-3)$, $B(5\,;1)$. Calculer les coordonnées de $\vec{AB}$ et sa norme.

\rubrique{Somme de vecteurs, relation de Chasles}
\exo{}\diff{1} Calculer $\vec u+\vec v$ avec $\vec u\begin{pmatrix}3\\-2\end{pmatrix}$ et
$\vec v\begin{pmatrix}1\\5\end{pmatrix}$.

\exo{}\diff{2} Calculer $\vec u+\vec v$ avec $\vec u\begin{pmatrix}-4\\6\end{pmatrix}$ et
$\vec v\begin{pmatrix}4\\-2\end{pmatrix}$.

\exo{}\diff{2} Sachant $\vec{AB}\begin{pmatrix}4\\2\end{pmatrix}$ et
$\vec{BC}\begin{pmatrix}3\\-5\end{pmatrix}$, calculer $\vec{AC}$ (relation de Chasles).

\exo{}\diff{1} Simplifier $\vec{AB}+\vec{BC}+\vec{CA}$.

\rubrique{Démontrer un parallélogramme}
\exo{}\diff{2} $A(0\,;0)$, $B(5\,;2)$, $C(8\,;6)$, $D(3\,;4)$. Montrer que $ABCD$ est un
parallélogramme.

\exo{}\diff{2} $A(1\,;1)$, $B(6\,;2)$, $C(8\,;5)$, $D(2\,;4)$. $ABCD$ est-il un
parallélogramme ? Justifier.

\exo{}\diff{2} $A(1\,;1)$, $B(6\,;2)$, $C(9\,;6)$, $D(4\,;5)$. Montrer que $ABCD$ est un
parallélogramme.

\exo{}\diff{3} $A(-2\,;0)$, $B(2\,;3)$, $D(-1\,;4)$. Déterminer les coordonnées du point
$C$ pour que $ABCD$ soit un parallélogramme.

% ═══════════════════════════════════════════════════════════════
\end{multicols}

\section{Exercices d'entraînement}
\begin{multicols}{2}

\rubrique{Vecteurs et translations}
\exo{}\diff{2} Calculer les coordonnées de l'image $A'$ du point $A(-3\,;5)$ par la
translation de vecteur $\vec u\begin{pmatrix}4\\-2\end{pmatrix}$.

\exo{}\diff{2} Le carré $ABCD$ a pour sommets $A(1\,;1)$, $B(4\,;1)$, $C(4\,;4)$,
$D(1\,;4)$. On le translate par $\vec u\begin{pmatrix}-2\\3\end{pmatrix}$. Donner les
coordonnées de l'image de chaque sommet.

\exo{}\diff{2} Une translation transforme $F(3\,;-2)$ en $F'(-1\,;6)$. Donner les
coordonnées du vecteur de cette translation.

\rubrique{Coordonnées et norme}
\exo{}\diff{2} $A(3\,;1)$, $B(7\,;4)$. Calculer les coordonnées de $\vec{AB}$ et sa norme.

\exo{}\diff{2} $A(-1\,;-2)$, $B(2\,;3)$. Calculer les coordonnées de $\vec{AB}$ et sa
norme (valeur exacte).

\exo{}\diff{2} $A(0\,;5)$, $B(6\,;-1)$. Calculer les coordonnées de $\vec{AB}$ et sa norme.

\exo{}\diff{2} $A(-3\,;4)$, $B(1\,;-2)$. Calculer les coordonnées de $\vec{AB}$ et sa
norme (valeur exacte).

\exo{}\diff{2} $A(2\,;2)$, $B(-2\,;-1)$. Calculer les coordonnées de $\vec{AB}$ et sa
norme.

\rubrique{Sommes et relation de Chasles}
\exo{}\diff{2} Calculer $\vec u+\vec v$ avec $\vec u\begin{pmatrix}-3\\4\end{pmatrix}$ et
$\vec v\begin{pmatrix}5\\-1\end{pmatrix}$.

\exo{}\diff{2} Sachant $\vec{AB}\begin{pmatrix}6\\-3\end{pmatrix}$ et
$\vec{AC}\begin{pmatrix}2\\5\end{pmatrix}$, calculer $\vec{BC}$ (indication :
$\vec{BC}=\vec{BA}+\vec{AC}$).

\exo{}\diff{2} Simplifier $\vec{MN}+\vec{NP}+\vec{PQ}$.

\rubrique{Parallélogrammes}
\exo{}\diff{2} $A(2\,;1)$, $B(6\,;3)$, $C(9\,;8)$, $D(5\,;6)$. Montrer que $ABCD$ est un
parallélogramme.

\exo{}\diff{2} $A(0\,;2)$, $B(4\,;0)$, $C(6\,;5)$, $D(2\,;6)$. $ABCD$ est-il un
parallélogramme ? Justifier.

\exo{}\diff{3} $A(1\,;2)$, $B(5\,;3)$, $C(4\,;-1)$. Déterminer les coordonnées du point
$D$ pour que $ABCD$ soit un parallélogramme.

\rubrique{Problèmes}
\exo{}\diff{2} Un randonneur part d'un point $A$ et se déplace du vecteur
$\vec u\begin{pmatrix}80\\60\end{pmatrix}$ (en mètres), puis du vecteur
$\vec v\begin{pmatrix}-30\\40\end{pmatrix}$. Calculer les coordonnées de sa position finale
par rapport à $A$, et sa distance totale à $A$ (à vol d'oiseau).

\exo{}\diff{2} Un drone décolle d'un point $O(0\,;0)$ et effectue le déplacement
$\vec u\begin{pmatrix}150\\0\end{pmatrix}$ puis $\vec v\begin{pmatrix}0\\200\end{pmatrix}$
(en mètres). Calculer sa distance à $O$ après ces deux déplacements.

% ═══════════════════════════════════════════════════════════════
\end{multicols}

\section{Sujets type BEPC}

\sujetbac[vecteurs, translation et triangle]
\begin{enumerate}[label=\arabic*)]
\item $A(1\,;2)$, $B(5\,;6)$. Calculer les coordonnées du vecteur $\vec{AB}$ et sa norme.
\item Calculer les coordonnées de l'image $A'$ du point $A(2\,;5)$ par la translation de
vecteur $\vec u\begin{pmatrix}3\\-4\end{pmatrix}$.
\item Calculer la somme des vecteurs $\vec u\begin{pmatrix}2\\-3\end{pmatrix}$ et
$\vec v\begin{pmatrix}-5\\7\end{pmatrix}$.
\item Le triangle $ABC$ a pour sommets $A(0\,;0)$, $B(4\,;1)$, $C(2\,;5)$. On le translate
par le vecteur $\vec u\begin{pmatrix}3\\-2\end{pmatrix}$. Donner les coordonnées de $A'$,
$B'$, $C'$.
\end{enumerate}

\sujetbac[Chasles et parallélogramme]
\begin{enumerate}[label=\arabic*)]
\item Sachant $\vec{AB}\begin{pmatrix}4\\3\end{pmatrix}$ et
$\vec{BC}\begin{pmatrix}1\\-5\end{pmatrix}$, calculer $\vec{AC}$.
\item $A(0\,;0)$, $B(5\,;2)$, $C(9\,;5)$, $D(4\,;3)$. Montrer que $ABCD$ est un
parallélogramme.
\item $A(-1\,;3)$, $B(4\,;15)$. Calculer la norme du vecteur $\vec{AB}$.
\item Simplifier $\vec{AB}+\vec{BC}+\vec{CD}$.
\end{enumerate}

\sujetbac[vecteurs et repérage d'un technicien]
\begin{enumerate}[label=\arabic*)]
\item $A(2\,;-1)$, $B(6\,;2)$. Calculer les coordonnées de $\vec{AB}$ et sa norme.
\item Le point $A(3\,;4)$ est translaté par le vecteur $\vec u\begin{pmatrix}-5\\2\end{pmatrix}$
pour donner $A'$. Calculer les coordonnées de $A'$.
\item $A(1\,;1)$, $B(4\,;3)$, $C(6\,;7)$, $D(3\,;5)$. $ABCD$ est-il un parallélogramme ?
Justifier.
\item Un technicien part d'une antenne repérée par le point $O(0\,;0)$ sur un plan. Il se
déplace du vecteur $\vec u\begin{pmatrix}120\\90\end{pmatrix}$ (en mètres), puis du vecteur
$\vec v\begin{pmatrix}-30\\30\end{pmatrix}$.
\begin{enumerate}[label=\alph*)]
\item Calculer les coordonnées de sa position finale.
\item Calculer sa distance à l'antenne $O$ (à vol d'oiseau).
\end{enumerate}
\end{enumerate}

% ═══════════════════════════════════════════════════════════════
\section{Corrigés}
\begin{multicols}{2}

\corrige{de l'exercice 1}
$\vec{AB}$ et $\vec{CD}$ ont les mêmes coordonnées, donc $\vec{AB}=\vec{CD}$ : ce sont des
représentants du même vecteur.

\corrige{de l'exercice 2}
$A'(2+3\,;5-4)=A'(5\,;1)$.

\corrige{de l'exercice 3}
$A'(3\,;-2)$, $B'(7\,;-1)$, $C'(5\,;3)$.

\corrige{de l'exercice 4}
$E=E'-\vec u$, soit $E(7-(-2)\,;-1-5)=E(9\,;-6)$.

\corrige{de l'exercice 5}
$\vec{AB}\begin{pmatrix}3\\4\end{pmatrix}$, $\|\vec{AB}\|=\sqrt{9+16}=5$.

\corrige{de l'exercice 6}
$\vec{AB}\begin{pmatrix}5\\12\end{pmatrix}$, $\|\vec{AB}\|=\sqrt{25+144}=13$.

\corrige{de l'exercice 7}
$\vec{AB}\begin{pmatrix}-6\\8\end{pmatrix}$, $\|\vec{AB}\|=\sqrt{36+64}=10$.

\corrige{de l'exercice 8}
$\vec{AB}\begin{pmatrix}3\\4\end{pmatrix}$, $\|\vec{AB}\|=\sqrt{9+16}=5$.

\corrige{de l'exercice 9}
$\vec u+\vec v\begin{pmatrix}4\\3\end{pmatrix}$.

\corrige{de l'exercice 10}
$\vec u+\vec v\begin{pmatrix}0\\4\end{pmatrix}$.

\corrige{de l'exercice 11}
$\vec{AC}\begin{pmatrix}7\\-3\end{pmatrix}$.

\corrige{de l'exercice 12}
$\vec{AB}+\vec{BC}+\vec{CA}=\vec{AA}=\vec0$.

\corrige{de l'exercice 13}
$\vec{AB}\begin{pmatrix}5\\2\end{pmatrix}$, $\vec{DC}\begin{pmatrix}5\\2\end{pmatrix}$.
Égaux : $ABCD$ est un parallélogramme.

\corrige{de l'exercice 14}
$\vec{AB}\begin{pmatrix}5\\1\end{pmatrix}$, $\vec{DC}\begin{pmatrix}6\\1\end{pmatrix}$.
Différents : $ABCD$ n'est \textbf{pas} un parallélogramme.

\corrige{de l'exercice 15}
$\vec{AB}\begin{pmatrix}5\\1\end{pmatrix}$, $\vec{DC}\begin{pmatrix}5\\1\end{pmatrix}$.
Égaux : $ABCD$ est un parallélogramme.

\corrige{de l'exercice 16}
$ABCD$ parallélogramme $\iff\vec{AB}=\vec{DC}$. $\vec{AB}\begin{pmatrix}4\\3\end{pmatrix}$,
donc $C=D+\vec{AB}=(-1+4\,;4+3)=(3\,;7)$.

\corrige{du sujet 1}
1) $\vec{AB}\begin{pmatrix}4\\4\end{pmatrix}$, $\|\vec{AB}\|=\sqrt{16+16}=4\sqrt2$.\\
2) $A'(5\,;1)$.\\
3) $\vec u+\vec v\begin{pmatrix}-3\\4\end{pmatrix}$.\\
4) $A'(3\,;-2)$, $B'(7\,;-1)$, $C'(5\,;3)$.

\corrige{du sujet 2}
1) $\vec{AC}\begin{pmatrix}5\\-2\end{pmatrix}$.\\
2) $\vec{AB}\begin{pmatrix}5\\2\end{pmatrix}$, $\vec{DC}\begin{pmatrix}5\\2\end{pmatrix}$.
Égaux : $ABCD$ est un parallélogramme.\\
3) $\|\vec{AB}\|=\sqrt{5^{2}+12^{2}}=\sqrt{169}=13$.\\
4) $\vec{AB}+\vec{BC}+\vec{CD}=\vec{AD}$.

\corrige{du sujet 3}
1) $\vec{AB}\begin{pmatrix}4\\3\end{pmatrix}$, $\|\vec{AB}\|=\sqrt{16+9}=5$.\\
2) $A'(3-5\,;4+2)=A'(-2\,;6)$.\\
3) $\vec{AB}\begin{pmatrix}3\\2\end{pmatrix}$, $\vec{DC}\begin{pmatrix}3\\2\end{pmatrix}$.
Égaux : $ABCD$ est un parallélogramme.\\
4a) Déplacement total $\begin{pmatrix}120-30\\90+30\end{pmatrix}=\begin{pmatrix}90\\120\end{pmatrix}$,
donc position finale $(90\,;120)$.\\
4b) Distance $=\sqrt{90^{2}+120^{2}}=\sqrt{22\,500}=150$ m.

\end{multicols}
\exercicesBanner
